\section{Related Work}
\label{sec:related_work}
% \subsection{Video Segmentation and 2D Editing}

% Video editing often relies on accurate object delineation and temporally consistent propagation. Classical video object segmentation (VOS) approaches propagate masks frame-by-frame using memory networks or tracking-based heads, enabling fine-grained, object-level edits and cut-and-paste workflows. Notable lines include semi-supervised VOS with space-time memory (STM)~\cite{oh2019videoobjectsegmentationusing}, long-range memory with efficient retrieval (XMem)~\cite{cheng2022xmemlongtermvideoobject}, and tracking-by-segmentation frameworks such as MaskTrack R-CNN~\cite{yang2019videoinstancesegmentation}. Complementary tools like video inpainting~\cite{li2022endtoendframeworkflowguidedvideo} and matting allow replacing or removing objects while maintaining background coherence. Although these methods deliver strong temporal consistency, they operate purely in 2D image space and do not enforce multi-view or 3D geometric consistency.


% Recent diffusion-based 2D image editors produce high-quality, instruction- or prompt-conditioned edits. Techniques include noise-to-image guidance and partial denoising (SDEdit)~\cite{meng2022sdeditguidedimagesynthesis}, cross-attention manipulation (Prompt-to-Prompt)~\cite{hertz2022prompttopromptimageeditingcross}, instruction tuning for localized edits (InstructPix2Pix)~\cite{ip2p}, and structure-conditional control (ControlNet)~\cite{controlnet}. These methods excel at preserving layout while modifying appearance or style, but naively applying them per frame (or per view) leads to flicker and identity drift without additional temporal or geometric constraints.

% To address temporal coherence, video extensions of diffusion models reuse or constrain cross-frame features. Representative methods include one-shot or few-shot tuning for videos (Tune-A-Video)~\cite{wu2023tuneavideooneshottuningimage}, zero-shot attention reuse (Text2Video-Zero)~\cite{khachatryan2023text2videozerotexttoimagediffusionmodels}, attention/token propagation (TokenFlow)~\cite{geyer2023tokenflowconsistentdiffusionfeatures}, edit-specific attention alignment (FateZero, Video-P2P)~\cite{qi2023fatezerofusingattentionszeroshot,liu2023videop2pvideoeditingcrossattention}, and layered/implicit canonicalization (CoDeF)~\cite{ouyang2024codefcontentdeformationfields}. Systems such as Rerender~\cite{yang2023rerendervideozeroshottextguided}, AnimateDiff~\cite{guo2024animatediffanimatepersonalizedtexttoimage}, and stable video diffusion models~\cite{blattmann2023stablevideodiffusionscaling} further strengthen temporal consistency via optical-flow guidance, cross-frame attention, or structural conditioning (depth/edges). Despite impressive short-range coherence, these approaches remain fundamentally 2D: they lack explicit 3D scene reasoning, so edits do not necessarily correspond to a single, consistent 3D configuration across viewpoints. When such 2D/video edits are lifted to supervision for 3D reconstruction, view-inconsistent artifacts or geometry drift can arise. This gap motivates 3D-aware editing frameworks that couple object semantics with geometry, ensuring both temporal and multi-view consistency.

% \subsection{3D Scene Understanding and Editing}
% Recent advancements in 3D neural reconstruction methods, notably Neural Radiance Fields (NeRFs) and 3D Gaussian Splatting (3DGS), have enabled photorealistic rendering from posed images. Building upon these techniques, researchers have sought to integrate semantic information to enhance the understanding and manipulation of reconstructed 3D scenes.
% %
% Many existing works leverage semantic features~\cite{feature3dgs} or 3D segmentation~\cite{gg, gaga} to achieve 3D scene editing of specified objects. NeRF-Art~\cite{nerfart} employs global-local contrastive learning, directional constraints in the CLIP embedding space, and weight regularization to ensure that NeRF preserves its original structure during text-guided stylization. Instruct-NeRF2NeRF~\cite{ip2p} reconstructs a 3D scene with NeRF~\cite{nerf} and then edits multi-view renderings with a 2D image editing method Instruct-Pix2Pix~\cite{ip2p}, in which a text-conditioned diffusion model is trained to edit the input image that aligns coherently with the input instruction prompt. The edited images are further utilized to continue optimizing the scene. GaussianCtrl~\cite{gaussctrl} employs ControlNet~\cite{controlnet} for image editing and achieves view consistency by depth-conditioned editing to enforce geometric consistency and attention-based latent code alignment to achieve appearance consistency. GaussianEditor~\cite{gaussianeditor} enhances precision and control in editing through a proposed Gaussian semantic tracing, which traces the editing target throughout the training process.

%To eliminate per-scene optimization, DN2N~\cite{} and Free-Editor~\cite{} utilize generalizable NeRF, however, 

% \subsection{Generalizable Feed-forward Reconstructors}

% Recent advancements in large 3D reconstruction models (LRMs)~\cite{mvsnerf, murf, pixelsplat, mvsplat, noposplat, gslrm, lgm, lsm, splatter_image, facelift} have significantly accelerated 3D reconstruction tasks. Leveraging scalable transformer architectures~\cite{transformer, vit} and extensive datasets like Objaverse~\cite{objaverse, objaversexl}, these models effectively learn generalizable 3D priors. In contrast to traditional per-scene optimization methods~\cite{nerf, gs, neus}, LRMs employ a feed-forward approach, enabling rapid inference of high-quality NeRF, mesh, or 3D Gaussian representations from sparse input images in less than a second. Specific approaches such as PixelSplat~\cite{pixelsplat} utilize epipolar geometry for depth estimation, whereas methods like MVSNeRF~\cite{mvsnerf}, MuRF~\cite{murf}, and MVSplat~\cite{mvsplat} construct cost volumes for aggregating multi-view information. MVSplat360~\cite{mvsplat360} extends MVSplat by enabling 360-degree feed-forward reconstruction refined through stable video diffusion models. Similarly, DepthSplat~\cite{depthsplat} incorporates explicit depth information to enhance novel-view synthesis. Approaches such as NoPoSplat~\cite{noposplat} leverage the pose estimation capabilities of MASt3R~\cite{dust3r, mast3r} to reconstruct scenes from unposed images, while LSM~\cite{lsm} integrates reconstruction with semantic scene understanding. Despite these advancements, the application of feed-forward reconstructors in 3D scene editing and manipulation remains relatively unexplored.
\subsection{2D Image and Video Editing}
Recent diffusion-based editors enable high-quality, instruction-conditioned edits through techniques such as noise-to-image guidance~\cite{meng2022sdeditguidedimagesynthesis}, cross-attention manipulation~\cite{hertz2022prompttopromptimageeditingcross}, instruction tuning~\cite{ip2p}, and structure-conditional control~\cite{controlnet}. While effective for single images, naively applying them per frame leads to flicker and identity drift.

To address temporal coherence, video extensions employ attention/token propagation~\cite{geyer2023tokenflowconsistentdiffusionfeatures}, edit-specific attention alignment~\cite{qi2023fatezerofusingattentionszeroshot,liu2023videop2pvideoeditingcrossattention}, optical-flow guidance~\cite{yang2023rerendervideozeroshottextguided}, and implicit canonicalization~\cite{ouyang2024codefcontentdeformationfields}. Despite strong short-range coherence, these approaches remain fundamentally 2D and lack explicit 3D scene reasoning, leading to view-inconsistent artifacts when lifted to 3D supervision.

\subsection{3D Scene Editing}
Neural reconstruction methods such as NeRF~\cite{nerf} and 3D Gaussian Splatting~\cite{3dgs} enable photorealistic rendering, motivating integration with semantic editing. Existing approaches typically follow a reconstruct-edit-refit pipeline. Instruct-NeRF2NeRF~\cite{in2n} reconstructs scenes with NeRF, edits multi-view renderings using InstructPix2Pix~\cite{ip2p}, and re-optimizes the scene with edited images. GaussianCtrl~\cite{gaussctrl} employs ControlNet~\cite{controlnet} with depth conditioning and attention-based alignment for view consistency, while GaussianEditor~\cite{gaussianeditor} introduces Gaussian semantic tracing for precise editing control. Though achieving impressive quality, these methods require slow per-scene optimization and are prone to multi-view inconsistencies.

\subsection{Generalizable Feed-forward Reconstructors}
Large reconstruction models (LRMs)~\cite{mvsnerf, pixelsplat, mvsplat, noposplat, gslrm, lgm, splatter_image, lsm} have revolutionized 3D reconstruction by learning generalizable priors from large-scale datasets like Objaverse~\cite{objaverse}. Leveraging transformer architectures~\cite{transformer, vit}, these models enable rapid inference of NeRF or Gaussian representations from sparse views without per-scene optimization. Methods such as PixelSplat~\cite{pixelsplat} utilize epipolar geometry, while many other methods~\cite{mvsnerf, murf, mvsplat, mvsplat360} construct cost volumes for multi-view aggregation. Despite their efficiency in reconstruction, the application of feed-forward reconstructors to 3D scene editing remains largely unexplored, which our work addresses.